% Generated by scripts/assemble.py; corrections belong in source/corrections.json.
\examhead{18}{中国科学院大学 2026 年数学分析试卷(回忆版)}{23}
% p024-b01; source PDF page 24
\blockanchor{p024-b01}
(注.题目来自考生回忆,部分表述或数据可能有误)\par
% p024-b02; source PDF page 24
\blockanchor{p024-b02}
\question{1}求极限.\par
% p024-b03; source PDF page 24
\blockanchor{p024-b03}
(1) $\lim_{n\to \infty}\frac{1}{n!}\sum_{p = 1}^{n}p!$.\par
% p024-b04; source PDF page 24
\blockanchor{p024-b04}
(2) $\lim_{\varepsilon \to 0^{+}}\varepsilon \ln \left(e^{\frac{a}{\varepsilon}} + e^{\frac{b}{\varepsilon}}\right).$\par
% p024-b05; source PDF page 24
\blockanchor{p024-b05}
\question{2}数列  $\{a_{n}\}$  满足  $a_{1}=1, a_{n+1}=\frac{1}{1+a_{n}}(n=1,2,\cdots)$ ，证明： $\{a_{n}\}$  收敛，并求其极限值.\par
% p024-b06; source PDF page 24
\blockanchor{p024-b06}
\question{3}求函数  $f(x_{1}, x_{2}, \cdots, x_{n}) = x_{1}^{2} + x_{2}^{2} + \cdots + x_{n}^{2}$  在条件  $\sum_{i=1}^{n} a_{i} x_{i} = 1$  下的最小值.\par
% p024-b07; source PDF page 24
\par\Needspace{6\baselineskip}
\blockanchor{p024-b07}
\question{4}函数 $f(x)$ 在 $[0, a]$ 上二阶可导，且 $|f''(x)| \leq M, x \in [0, a]$. 若 $f(x)$ 在 $(0, a)$ 上存在最大值，证明：\par
\nopagebreak[4]
% p024-b08; source PDF page 24
\blockanchor{p024-b08}
\displayformula{\left| f ^ {\prime} (0) \right| + \left| f ^ {\prime} (a) \right| \leq M a.}
% p024-b09; source PDF page 24
\par\Needspace{6\baselineskip}
\blockanchor{p024-b09}
\question{5}求曲面积分\par
\nopagebreak[4]
% p024-b10; source PDF page 24
\blockanchor{p024-b10}
\displayformula{\iint_ {\Sigma} x ^ {3} \mathrm{d} y \mathrm{d} z + y ^ {3} \mathrm{d} z \mathrm{d} x + z ^ {3} \mathrm{d} x \mathrm{d} y.}
% p024-b11; source PDF page 24
\blockanchor{p024-b11}
其中 $\Sigma$ 为上半椭球面 $\frac{x^2}{a^2} +\frac{y^2}{b^2} +\frac{z^2}{c^2} = 1(z\geq 0)$ ，取上侧.\par
% p024-b12; source PDF page 24
\par\Needspace{6\baselineskip}
\blockanchor{p024-b12}
\question{6}设 $P(x, y)$, $Q(x, y)$ 存在连续的偏导数, 且对任意以 $(x_0, y_0) \in \mathbb{R}^2$ 为圆心以 $r > 0$ 为半径的上半圆周 $L: x = x_0 + r\cos \theta, y = y_0 + r\sin \theta, \theta \in [0, \pi]$, 都有\par
\nopagebreak[4]
% p024-b13; source PDF page 24
\blockanchor{p024-b13}
\displayformula{\int_ {L} P (x, y) \mathrm{d} x + Q (x, y) \mathrm{d} y = 0.}
% p024-b14; source PDF page 24
\blockanchor{p024-b14}
证明 $P(x,y) = \frac{\partial Q}{\partial x} (x,y)\equiv 0.$\par
% p024-b15; source PDF page 24
\blockanchor{p024-b15}
\question{7}求积分  $\int_{0}^{+\infty}\frac{\cos(ax)}{1+x^{2}}dx.$\par
% p024-b16; source PDF page 24
\par\Needspace{6\baselineskip}
\blockanchor{p024-b16}
\question{8}已知 $F_{n} = \int \left(\frac{\sin\frac{x - a}{2}}{\sin\frac{x + a}{2}}\right)^{n}\mathrm{d}x,$ 其中 $0 < a < \pi$ ，证明：\par
\nopagebreak[4]
% p024-b17; source PDF page 24
\blockanchor{p024-b17}
\displayformula{F _ {n} = \frac {2 \sin a}{n - 1} \left(\frac {\sin \frac {x - a}{2}}{\sin \frac {x + a}{2}}\right) ^ {n - 1} + 2 \cos a F _ {n - 1} - F _ {n - 2}.}
% p024-b18; source PDF page 24
\par\Needspace{6\baselineskip}
\blockanchor{p024-b18}
\question{9}证明：\par
\nopagebreak[4]
% p024-b19; source PDF page 24
\blockanchor{p024-b19}
\displayformula{\iint_ {D} f (m \sin \theta \cos \varphi + n \sin \theta \sin \varphi + p \cos \theta) \sin \theta \mathrm{d} \theta \mathrm{d} \varphi = 2 \pi \int_ {- 1} ^ {1} f (u \sqrt {m ^ {2} + n ^ {2} + p ^ {2}}) \mathrm{d} u.}
% p024-b20; source PDF page 24
\blockanchor{p024-b20}
其中 $D: \theta \in [0, \pi]$, $\varphi \in [0, 2\pi]$.\par
% p024-b21; source PDF page 24
\blockanchor{p024-b21}
\question{10}证明: $\lim_{n\to \infty}\int_{a}^{b}f(x)|\sin nx|\mathrm{d}x = \frac{2}{\pi}\int_{a}^{b}f(x)\mathrm{d}x.$\par
